\section{The smallest service that works}
\label{sec:ch02-smallest-service}
\index{HotChocolate!getting started|(}

\index{HotChocolate!templates}%
The fastest route into HotChocolate~16 is not a package reference. It is a
template, which ChilliCream ships on NuGet:

\begin{minted}{text}
dotnet new install HotChocolate.Templates
dotnet new graphql -o src/Mosaic.Api
\end{minted}

Note the package name. ChilliCream is the company, so prefixing it with
\code{ChilliCream.} is the natural guess, and no such package
exists~\autocite{nuget2026templates}. Installing the real one registers three
templates: \code{graphql} for a server, \code{graphql-gateway} for a Fusion
gateway, and \code{graphql-azf} for an Azure Function. We want the first.

I recommend starting from the template rather than assembling the project by
hand, and not for the usual reason. It is not that the setup is hard. It is that
HotChocolate~16 changed how a server is wired up, and almost everything written
about HotChocolate on the public internet describes an older version. The
template ships from ChilliCream at the same version as the packages, so it
matches the release you are installing. Here is what it produces, in full:

\begin{minted}{csharp}
var builder = WebApplication.CreateBuilder(args);

builder.AddGraphQL().AddTypes();

var app = builder.Build();

app.MapGraphQL();

app.RunWithGraphQLCommands(args);
\end{minted}

\index{HotChocolate!AddGraphQL}%
If you have used HotChocolate before, the second line is the surprising one. The
call you probably remember is this:

\begin{minted}{csharp}
builder.Services
    .AddGraphQLServer()
    .AddQueryType<Query>();
\end{minted}

That still compiles. It is not deprecated, and nothing in this book will break
if you keep writing it. But it is no longer the shape the documentation teaches,
and in 16.6.0 there are three similarly named entry points with genuinely
different behaviour~\autocite{chillicream2026source}.
\code{builder.AddGraphQL()}, on the host builder, forwards to
\code{AddGraphQLServer()} and applies the ASP.NET Core defaults.
\code{builder.Services.AddGraphQL()}, on the service collection, does not apply
them. One character of difference in what you type, and a materially different
server at the other end, so it is worth knowing before you need it.

\index{HotChocolate!source generator}%
The more interesting absence is \code{Query}. Nothing in that file names a type,
and yet the service serves a schema. \code{AddTypes()} is generated at compile
time by a Roslyn source generator that ships in the
\code{HotChocolate.Types.Analyzers} package, which the template references as an
analyzer. The generator scans the assembly for types carrying attributes it
recognises, and writes the registration code that the file above then calls.
The template's query type looks like this:

\begin{minted}{csharp}
[QueryType]
public static partial class Query
{
    public static Book GetBook()
        => new Book("C# in depth.", new Author("Jon Skeet"));
}
\end{minted}

A static class, marked \code{partial} so the generator can add to it, with
static methods that become fields. \code{GetBook} becomes \code{book}:
HotChocolate strips the \code{Get} prefix, and it strips an \code{Async} suffix
too, so \code{GetProductsAsync} becomes \code{products}.

\index{HotChocolate!Module attribute}%
One detail here bites early and is worth showing before it does. The generated
method is not always called \code{AddTypes}. Its name comes from an
assembly-level attribute, which the template ships in a file of its own as
\code{Module("Types")}. Rename the module to something meaningful for your
service and the generated method is renamed with it. Mosaic uses this:

\begin{minted}{csharp}
[assembly: Module("Mosaic")]
\end{minted}

so its registration call is \code{AddMosaic()}. If you rename the module and
forget to update \code{Program.cs} to match, you do not get a helpful message
about a missing method. \code{AddTypes} still resolves, to a pair of unrelated
\code{params} overloads that take arrays of types, and the compiler reports:

\begin{minted}{text}
error CS0121: The call is ambiguous between the following methods or
properties: 'SchemaRequestExecutorBuilderExtensions.AddTypes(
IRequestExecutorBuilder, params System.Type[])' and
'SchemaRequestExecutorBuilderExtensions.AddTypes(
IRequestExecutorBuilder, params HotChocolate.Types.ITypeDefinition[])'
\end{minted}

Which is an accurate description of the compiler's problem and no help at all
with yours. The attribute is assembly-scoped rather than folder-scoped, which
matters more than it looks: one \code{AddMosaic()} call registers every
GraphQL type in the assembly, no matter which folder it sits in. That single
fact is what makes the layout in section~\ref{sec:ch02-six-domains} possible.

\index{Nitro}%
Run the project and open \code{http://localhost:5100/graphql} in a browser.
HotChocolate notices that a browser asked and serves Nitro, its own IDE, from
the same endpoint that answers queries~\autocite{chillicream2026endpoints}.
Nitro was called Banana Cake Pop until October 2024, when Rafael Staib announced
the rename on ChilliCream's blog~\autocite{staib2024nitro}. The old name has been
scrubbed from the current documentation, so if you find a page using it you are
reading something out of date. The same endpoint also answers \code{?sdl} with
the schema as a file, which is occasionally more useful than the IDE.

Point a client at it, ask for the one field that exists, and you have a working
GraphQL server. That is the \emph{quickly} part, and it is over. Everything
after this is about the choices the template made silently on your behalf.
\index{HotChocolate!getting started|)}
